\section{Lagerarkitekturen}
\label{sec:arch}

\bookfigure{architecture}{Drivrutinsstacken. Varje vågrät linje är en gräns någon inte får titta
igenom.}{fig:arch}

Två sömmar bär hela konstruktionen, och de löser olika problem:

\begin{booktable}{@{}lL@{}}
\thead{Söm} & \thead{Gör vad testbart} \\ \midrule
\code{driver::can::Interface} & \textbf{Applikationen.} \code{EchoNode} kan köras mot en
  \code{Stub} och behöver aldrig veta att hårdvara finns. \\
\code{driver::can::ByteTransport} & \textbf{Drivrutinen själv.} Samma \code{Spi}-kod som flashas
  till målet kan köras på din laptop mot en transport du skriptar. \\
\end{booktable}

Den andra sömmen är inte uppenbar, och den kommer i \kapref{ch:byte}. Utan den går \code{Spi} bara
att köra på riktig hårdvara --- och den finns först i \kapref{ch:bringup}, näst sist i kursen, på
ett kort som samtidigt felsöks av någon annan.

\textbf{Regeln som följer, och som en granskare ska leta efter: ingen fil ovanför
\code{driver/can/interface.hpp} får innehålla ordet SPI.}

\subsection{Sömmen har redan lönat sig en gång}

Argumentet för ett interface brukar vara hypotetiskt --- *tänk om hårdvaran byts ut*. Här är det
inte hypotetiskt.

En tidigare version av den här konstruktionen nådde registren via \textbf{minnesmappad I/O} på
basadressen \code{0xFF200000}. Det förutsätter en processor med en adressbuss ut mot FPGA-väven,
och mikrokontrollern i det här systemet har ingen sådan buss. Registren nås över SPI i stället.

Allt ovanför \code{Interface} hade överlevt den ändringen \textbf{oförändrat}. Det är hela
argumentet, demonstrerat på riktig kod.

\subsection{Sändningssekvensen}

Att skicka en frame (\code{id 0x123}, \code{dlc 2}, data \code{AA BB}) är fem registerskrivningar
och en pollning:

\begin{console}
1. läs   STATUS        -> bit 0 satt? annars: upptaget, returnera false
2. skriv TX_ID         = 0x00000123
3. skriv TX_DLC        = 0x00000002
4. skriv TX_DATA_LO    = 0xAABB0000
5. skriv TX_DATA_HI    = 0x00000000
6. skriv TX_SEND       = 0x00000001      -> STATUS bit 0 går låg
   ...
7. polla STATUS        -> bit 0 tillbaka: porten är fri
8. läs   ERROR_FLAGS   -> 0 = framen gick ut, 1 = något gick fel
\end{console}

Och att ta emot en:

\begin{console}
1. läs   STATUS        -> bit 1 satt? annars: ingen frame, returnera false
2. läs   RX_ID, RX_DLC
3. läs   RX_DATA_LO, RX_DATA_HI, packa upp, maskera mot RX_DLC
4. skriv RX_ACK        = 0x00000001      -> STATUS bit 1 går låg
\end{console}

\code{TX_SEND} \textbf{måste} skrivas sist. Den skrivningen betyder inte "spara ett värde" utan
"börja sända nu", och kontrollern läser då de övriga registren. Skriv triggern först och du
skickar förra framens innehåll --- vilket vid bring-upen ser ut som en nod som ligger exakt en
frame efter.
